\documentclass[10pt, letterpaper]{article}

% --- PACKAGES AND SETUP ---
\usepackage[utf8]{inputenc}
\usepackage[T1]{fontenc}
\usepackage[english]{babel}
\usepackage[margin=1in]{geometry}
\usepackage{parskip}
\usepackage{amsmath}
\usepackage{amssymb}
\usepackage{soul}
\usepackage{csquotes}
\usepackage{ragged2e}
\usepackage[none]{hyphenat}
\usepackage{chngcntr}
\usepackage{enumitem}
\usepackage{caption}

% --- BIBLIOGRAPHY SETUP (biblatex) ---
\usepackage[
    backend=biber,
    style=authoryear,
    sorting=nyt,
    maxcitenames=2,
    mincitenames=1 
]{biblatex}

\addbibresource{grass.bib}
\counterwithin{subsection}{section}

% --- HYPERLINKS AND PDF METADATA ---
\usepackage{hyperref}
\hypersetup{
    colorlinks=true,
    linkcolor=blue,
    citecolor=blue,
    urlcolor=blue,
    pdftitle={Thesis Layout: Golf Course Opportunity Costs},
    pdfauthor={Michael Manolakis},
}

% --- DOCUMENT START ---
\begin{document}

\begin{center}
    {\huge Analysis: Green Golf Courses and Money Funding \\}
    {\large The Aggregate Opportunity Cost of Recreational Land Use: A Counterfactual Valuation of U.S. Golf Courses \\}
    {\normalsize Michael Manolakis \textbf{\textemdash} April 1, 2026 \\}
\end{center}

\vspace{1cm}

\begin{abstract}
This thesis investigates the macroeconomic inefficiency of recreational land allocation in the United States, specifically quantifying the opportunity cost of maintaining golf courses in supply-constrained markets. Utilizing a novel spatial dataset of 16,297 facilities, this study applies the real estate appraisal principle of Highest and Best Use (HBU) to establish residential and agricultural counterfactuals. The methodology integrates Federal Housing Finance Agency (FHFA) and United States Department of Agriculture (USDA) land valuations, employing a rural-urban continuum bifurcation to control for spatial elasticity. To rigorously resolve 38\% missing acreage data, the analysis utilizes Multiple Imputation by Chained Equations (MICE) and pools the parameters using Rubin's Rules. Empirical results reveal a pooled aggregate opportunity cost of approximately \$244.9 billion. This locked land value is characterized by extreme spatial heterogeneity; while urban courses comprise 72\% of the total recreational acreage, they account for 98.5\% of the aggregate financial opportunity cost, reflecting the massive regulatory ``zoning tax'' endemic to metropolitan areas. Despite this immense latent value, the market fails to transition these parcels to their HBU due to severe institutional frictions. A dual layer of exclusionary municipal zoning and inflexible private restrictive covenants effectively paralyzes land-use transitions, rendering Coasian bargaining impossible. These findings highlight a systemic market failure, suggesting that unlocking this land for high-density housing or utility-scale renewable energy infrastructure requires aggressive legal and regulatory reform.
\end{abstract}

\section{Introduction}

The allocation and regulation of urban and suburban land represent critical challenges within modern spatial economics. In many major United States metropolitan areas, the supply of developable land is highly inelastic, driven largely by restrictive land-use regulations rather than physical geographic constraints \parencite{Glaeser2017, Gyourko2021}. This artificial scarcity has led to a pronounced divergence between the physical cost of construction and the market price of housing, generating a substantial ``zoning tax'' that inflates residential land values \parencite{Glaeser2005, Davis2021}. 

Amidst this intense competition for spatial resources, the American golf course represents a massive and highly specific land-use typology. Historically, real estate developers utilized golf courses as ``surrogate public goods'' to cross-subsidize the construction of master-planned communities, capturing exorbitant price premiums on adjacent residential lots \parencite{Limehouse2012}. However, empirical evidence suggests that the marginal amenity value of proximity to golf courses has depreciated over time as consumer preferences shift toward natural open spaces \parencite{Cho2009}. Despite this structural decline in recreational demand and the escalating value of residential land, hundreds of thousands of acres of prime real estate remain locked in this single-use capacity \parencite{Weinand2025}.

To quantify the macroeconomic inefficiency of this spatial allocation, this thesis applies the foundational real estate appraisal principle of Highest and Best Use (HBU). HBU dictates that land should be valued based on its most profitable, legally permissible, and physically possible use \parencite{AppraisalInstitute2020}. By treating the existing golf course as a baseline, we can establish a counterfactual valuation model to determine the aggregate opportunity cost ($V_{OC}$) of this land. The core economic framework is defined by the following equation:

\begin{equation}
    V_{OC} = \max(V_{Res}, V_{Ag}) - V_{Current}
\end{equation}

where $V_{Res}$ represents the land value under a residential development counterfactual, $V_{Ag}$ represents the value under an agricultural counterfactual, and $V_{Current}$ represents the capitalized value of the land in its current recreational use. 

In a frictionless, perfectly competitive market, land would naturally transition to the use that maximizes its residual value. However, the transition of golf course land to its HBU is severely impeded by a dual layer of institutional frictions. Publicly, municipal zoning ordinances frequently prohibit high-density redevelopment. Privately, properties are often encumbered by stale restrictive covenants and Homeowner Association (HOA) agreements that legally mandate the preservation of the golf course in perpetuity, regardless of its financial feasibility \parencite{Ellickson2022, Korngold2024}. 

Because these legal barriers prevent fluid Coasian bargaining and market optimization, this thesis employs a formal \textit{Hypothetical Condition} \parencite{AppraisalInstitute2020}. The analysis proceeds under the assumption of a frictionless legal environment \textemdash specifically, the hypothetical dissolution of both exclusionary zoning restrictions and parallel private covenants. This methodological assumption is necessary to isolate and measure the pure economic opportunity cost of the land.

Therefore, this thesis seeks to answer the following core research question: 

\begin{quote}
\textit{What is the aggregate financial opportunity cost of U.S. golf courses when evaluated against their Highest and Best Use counterfactuals of residential, urban, or agricultural development?}
\end{quote}
\textit{}

To answer this question, this study constructs a national-level spatial dataset integrating golf course boundaries from OpenStreetMap (OSM) and the Tigris Census database, paired with granular land valuation data from the Federal Housing Finance Agency (FHFA) and the United States Department of Agriculture (USDA). To address missing spatial covariates within the dataset, the analysis utilizes Multiple Imputation by Chained Equations (MICE), ensuring that the final aggregate valuation remains statistically robust and accounts for between-imputation variance \parencite{Burren2018, Rubin1987}. Ultimately, by quantifying this locked land value, this research aims to inform broader policy discussions regarding zoning reform, housing supply, and the economic efficiency of urban land use.


\subsection{Roadmap}
The remainder of this thesis is organized as follows: 

\textbf{Section 2} reviews the relevant economic and legal literature, establishing the macroeconomic context of regulatory land scarcity, the microeconomic pricing of recreational amenities, and the institutional frictions that govern land-use transitions. 

\textbf{Section 3} formalizes the theoretical framework, translating the appraisal doctrine of Highest and Best Use (HBU) into a neoclassical spatial model to mathematically define the aggregate opportunity cost function ($V_{OC}$). 

\textbf{Section 4} details the data and methodology, outlining the integration of geospatial boundaries with federal land valuation datasets (FHFA and USDA). This section also explicitly defines the hybrid valuation algorithm and the Multiple Imputation by Chained Equations (MICE) protocol utilized to resolve missing spatial covariates. 

\textbf{Section 5} presents the empirical results, reporting the pooled aggregate valuation of \$244.9 billion, decomposing this opportunity cost across urban and rural submarkets, and providing the pooled logarithmic regression outputs. 

\textbf{Section 6} discusses the macroeconomic and policy implications of these findings, analyzing the failure of Coasian bargaining due to zoning and private covenants, exploring renewable energy counterfactuals, and addressing methodological limitations regarding bulk-parcel valuation. 

Finally, \textbf{Section 7} concludes the thesis with a summary of the empirical findings and directions for future economic research.

\newpage

\section{Literature Review}


This thesis bridges three distinct strands of economic and legal literature: the macroeconomic valuation of urban land and regulatory scarcity, the microeconomic pricing of recreational amenities, and the institutional frictions that govern land-use transitions. By synthesizing these domains, this research establishes the theoretical foundation necessary to quantify the aggregate opportunity cost of golf course land.

\subsection{Urban Land Valuation and Regulatory Scarcity}
The primary driver of the opportunity cost associated with underutilized urban land is the escalating value of residential real estate ($V_{Res}$). Foundational work by \textcite{Davis2008} decomposes housing prices into structure costs and land values, demonstrating that the vast heterogeneity in U.S. housing prices is driven almost entirely by the price of the underlying land, rather than physical construction costs. \textcite{Davis2021} expand upon this by providing highly granular, tract-level estimates of residential land prices, confirming that land values follow steep spatial gradients that peak sharply near central business districts. 

This extreme appreciation in land value is largely attributed to artificial supply constraints. \textcite{Glaeser2005} introduce the concept of the ``zoning tax,'' defined as the gap between the market price of a housing unit and its marginal cost of production. In a frictionless market, free entry ensures that prices equal average costs; however, restrictive land-use regulations sever this relationship, bidding up the extensive margin value of land \parencite{Gyourko2021}. The macroeconomic implications of this inelastic housing supply are severe. \textcite{Glaeser2017} argue that regulatory constraints in high-productivity markets prevent the efficient spatial allocation of labor, resulting in significant deadweight loss to aggregate national output. Consequently, the opportunity cost of maintaining large tracts of single-use recreational land in these constrained markets is not merely a localized inefficiency, but a macroeconomic distortion.

\subsection{The Microeconomics of Recreational Amenities}
To accurately model the current value of golf courses ($V_{Current}$), it is necessary to understand the economic mechanisms that generate their revenue. Utilizing the two-stage hedonic pricing model established by \textcite{Rosen1974}, \textcite{Limehouse2010} demonstrate that environmental quality and recreational exclusivity are credence goods capitalized directly into golf course pricing. 

However, the primary economic driver of golf course development is often not the recreational operation itself, but the ``joint production'' of the course and adjacent real estate \parencite{Limehouse2012}. Developers utilize the golf course as a surrogate public good to cross-subsidize the project, extracting massive price premiums from the surrounding residential lots. While this model historically maximized developer profit, longitudinal hedonic analyses indicate that the amenity value of golf courses is subject to structural depreciation. \textcite{Cho2009} find that while proximity to open space acts as a substitute for private lot size, consumer preferences have shifted significantly over time. The marginal value of proximity to natural amenities (e.g., public parks, water bodies) has increased, whereas the premium associated specifically with golf courses has experienced a measurable decline. Furthermore, emerging literature by \textcite{Weinand2025} highlights the severe spatial inefficiency of golf courses, noting that the land footprint of these facilities vastly exceeds the land required for utility-scale renewable energy infrastructure, introducing alternative, non-residential counterfactuals for land optimization.

\subsection{Institutional Frictions and Land Use Rigidity}
If the residential value of land ($V_{Res}$) exceeds its current recreational value ($V_{Current}$), standard neoclassical theory dictates that the land should be redeveloped. The failure of the market to execute this transition is explained by severe institutional frictions and transaction costs, which prevent fluid Coasian bargaining \parencite{Coase1960}. 

The rigidity of current land use is enforced by a dual layer of public and private regulation. Publicly, municipal zoning laws frequently prohibit the densification required to achieve Highest and Best Use \parencite{Fischel2015}. While these planning policies generate localized environmental amenities, \textcite{Cheshire2002} demonstrate that the net welfare effect of such restrictive regulatory regimes is substantially negative and acts as a regressive tax on housing affordability. 

Privately, land-use transitions are blocked by restrictive covenants. \textcite{Ellickson2022} notes that the majority of modern housing developments are tethered by covenants that mandate the preservation of existing land uses (such as golf courses) in perpetuity. Because these covenants typically require supermajority votes to amend, they suffer from severe collective action problems and status quo bias, effectively freezing the land use. To resolve this, \textcite{Korngold2024} argues that legislative zoning reform is insufficient unless parallel private covenants are simultaneously voided under public policy doctrines. This literature confirms that the legal barriers to redevelopment are formidable, thereby justifying this study's use of a Hypothetical Condition to calculate the pure economic opportunity cost absent these specific legal frictions.

\subsection{Synthesis and Contribution}
While existing literature extensively documents the regulatory inflation of residential land values and the microeconomic pricing of golf course amenities, there is a distinct lack of empirical research quantifying the aggregate spatial intersection of these two phenomena. This thesis contributes to the literature by applying rigorous real estate appraisal frameworks \parencite{AppraisalInstitute2020} to granular spatial datasets, calculating the precise national opportunity cost of maintaining depreciating recreational assets in supply-constrained housing markets.


\newpage

\section{Theoretical Framework}

This section establishes the formal economic model utilized to quantify the spatial misallocation of recreational land. The framework integrates neoclassical urban economics with the professional appraisal doctrine of Highest and Best Use (HBU) to derive a mathematical function for aggregate opportunity cost.

\subsection{Neoclassical Land Allocation and Highest and Best Use}
In a frictionless, perfectly competitive spatial market, land is allocated to the use that generates the maximum residual value. This principle aligns with the appraisal standard of Highest and Best Use (HBU), defined as the reasonably probable and legal use of vacant land or an improved property that is physically possible, appropriately supported, financially feasible, and results in the highest value \parencite{AppraisalInstitute2020}. 

For a given parcel of land $i$, let $V_{i,k}$ represent the present value of the land under use $k$, where $k \in \{1, 2, ..., K\}$ represents the set of all physically possible and financially feasible uses. Under conditions of spatial equilibrium, the market value of the parcel, $V_i^*$, is determined by the use that maximizes this value:

\begin{equation}
    V_i^* = \max \{V_{i,1}, V_{i,2}, ..., V_{i,K}\}
\end{equation}

When the current use of the land, $V_{Current}$, diverges from $V_i^*$, an economic inefficiency exists. The magnitude of this inefficiency is the opportunity cost.

\subsection{The Opportunity Cost Function}
To quantify the inefficiency of land currently allocated to golf courses, this analysis restricts the set of counterfactual uses $K$ to the two most probable alternative land uses based on spatial geography: residential development ($V_{Res}$) and agricultural production ($V_{Ag}$). 

The fundamental opportunity cost ($V_{OC}$) for a specific golf course parcel $i$ is defined as the difference between its value under its HBU counterfactual and its value in its current recreational state:

\begin{equation}
    V_{OC, i} = \max(V_{Res, i}, V_{Ag, i}) - V_{Current, i}
\end{equation}

The aggregate national opportunity cost is the summation of $V_{OC, i}$ across all $N$ golf courses in the dataset:

\begin{equation}
    Aggregate~V_{OC} = \sum_{i=1}^{N} \left[ \max(V_{Res, i}, V_{Ag, i}) - V_{Current, i} \right]
\end{equation}

\subsubsection{Deconstructing the Counterfactuals}
The variables within the core equation are defined by distinct microeconomic determinants:

\begin{itemize}
    \item $\mathbf{V_{Res}}$ \textbf{(Residential Value):} The extensive margin value of the land if subdivided and developed for single-family or multi-family housing. In highly regulated metropolitan statistical areas (MSAs), $V_{Res}$ is artificially inflated by a regulatory ``zoning tax'' ($\tau_z$), which represents the scarcity premium induced by supply-side restrictions \parencite{Glaeser2005, Gyourko2021}. Therefore, $V_{Res}$ is a function of organic housing demand ($D_h$), physical construction costs ($C_c$), and regulatory scarcity ($\tau_z$).
    \item $\mathbf{V_{Ag}}$ \textbf{(Agricultural Value):} The capitalized value of the land's agricultural yield, primarily applicable to courses located in rural or exurban geographies where residential demand ($D_h$) is insufficient to justify development costs.
    \item $\mathbf{V_{Current}}$ \textbf{(Current Recreational Value):} The capitalized net operating income of the golf course. This value is heavily dependent on the joint production model, where the course serves as a surrogate public good that generates price premiums for adjacent real estate \parencite{Limehouse2012}. 
\end{itemize}

\subsection{Institutional Frictions and Market Failure}
If $V_{OC, i} > 0$, rational economic actors should execute a transaction to redevelop the land. The persistence of golf courses in high-value urban areas indicates a market failure driven by severe institutional frictions ($F$). 

These frictions consist of public regulatory constraints ($Z$), such as municipal zoning ordinances that prohibit high-density development, and private contractual constraints ($C$), such as restrictive covenants enforced by Homeowner Associations (HOAs) \parencite{Ellickson2022, Korngold2024}. Furthermore, the redevelopment of a golf course imposes a negative externality on adjacent homeowners by destroying their capitalized open-space amenity, leading to intense political resistance and prohibitive Coasian transaction costs \parencite{Coase1960}.

Mathematically, redevelopment only occurs if the opportunity cost exceeds the total institutional frictions:

\begin{equation}
    V_{OC, i} > F(Z_i, C_i)
\end{equation}

Because $F(Z_i, C_i)$ approaches infinity in many master-planned communities due to the legal impossibility of breaking stale covenants, the market cannot clear. To isolate the pure economic value of the land, this thesis applies a \textit{Hypothetical Condition}, setting $F(Z_i, C_i) = 0$. This allows for the unconstrained calculation of $V_{OC}$ based strictly on the economic parameters of $V_{Res}$ and $V_{Ag}$.

\newpage

\section{Data and Methodology}

To quantify the aggregate opportunity cost of U.S. golf courses, this study constructs a novel, national-level spatial dataset. The methodology integrates geospatial boundary manipulation, a hybrid valuation algorithm based on rural-urban continuum classifications, and advanced imputation techniques to resolve missing covariates.

\subsection{Dataset Scope and Sourcing}
The primary dataset consists of $N = 16,297$ golf course observations across the United States. Spatial point coordinates (latitude and longitude) for each facility, mapped to the WGS84 coordinate reference system, were spatially joined to U.S. Census Bureau County shapefiles (2022 vintage, 20-meter resolution) utilizing the \texttt{tigris} package in R. 

To establish the Highest and Best Use (HBU) counterfactuals, two distinct land valuation datasets were integrated:
\begin{enumerate}
    \item \textbf{Residential Land Values ($V_{Res}$):} Sourced from the Federal Housing Finance Agency (FHFA) 2022 dataset at the county level, joined via 5-digit Federal Information Processing Standard (FIPS) codes.
    \item \textbf{Agricultural Land Values ($V_{Ag}$):} Sourced from the United States Department of Agriculture (USDA) 2025 projections at the state level.
\end{enumerate}

\subsection{Data Dictionary}
The empirical models utilize the core variables defined in Table \ref{tab:data_dict}.

\begin{table}[h!]
\centering
\caption{Core Variables and Data Dictionary}
\label{tab:data_dict}
\renewcommand{\arraystretch}{1.2}
\begin{tabular}{p{3.5cm} p{2cm} p{7cm} p{2.5cm}}
\hline
\textbf{Variable Name} & \textbf{Data Type} & \textbf{Description} & \textbf{Unit} \\
\hline
\texttt{course\_id} & Numeric & Unique identifier for each observation. & Count \\
\texttt{final\_acreage} & Numeric & Total footprint of the facility (observed or imputed). & Acres \\
\texttt{holes} & Numeric & Total number of holes at the facility. & Count \\
\texttt{course\_type} & Categorical & Operational classification (e.g., Public, Private). & Text \\
\texttt{county\_fips} & Categorical & 5-digit Federal Information Processing Standard code. & Text \\
\texttt{RUCC\_2013} & Categorical & USDA Rural-Urban Continuum Code (1--9). & Index \\
\texttt{county\_type} & Categorical & Derived binary geographic classification (Urban/Rural). & Text \\
\texttt{Estimated\_Price} & Numeric & Proxied land value per acre assigned via hybrid model. & USD \\
\texttt{estimated\_value} & Numeric & Calculated total land value for the respective course. & USD \\
\hline
\end{tabular}
\end{table}

\subsection{The Hybrid Valuation Algorithm}
Because land supply elasticity and HBU vary drastically across geographic submarkets \parencite{Davis2008}, applying a uniform residential counterfactual nationally would systematically overestimate opportunity costs in inelastic regions. To resolve this, a hybrid valuation algorithm was developed utilizing the USDA Rural-Urban Continuum Codes (RUCC).

Observations were bifurcated based on their RUCC designation:
\begin{itemize}
    \item \textbf{Urban Classification (RUCC 1, 2, 3):} Courses located in metropolitan counties. The HBU counterfactual is assumed to be residential development. These observations were assigned the FHFA county-level residential price per acre ($V_{Res}$).
    \item \textbf{Rural Classification (RUCC 4 through 9):} Courses located in non-metropolitan or rural counties. The HBU counterfactual is assumed to be agricultural production. These observations were assigned the USDA state-level agricultural price per acre ($V_{Ag}$).
\end{itemize}
To ensure dataset completeness, a fallback protocol was implemented: any urban-classified geometry missing FHFA metrics was forcefully defaulted to the USDA agricultural state-level parameter, ensuring a conservative lower-bound valuation. The baseline opportunity cost for each parcel was then calculated as:
\begin{equation}
    Estimated\_Land\_Value_i = final\_acreage_i \times Estimated\_Price\_Per\_Acre_i
\end{equation}

\subsection{Handling Missing Covariates: Multiple Imputation}
A critical limitation of the primary dataset was the absence of acreage data for exactly 38\% of the golf course observations. Basic listwise deletion of these rows would introduce severe selection bias and artificially deflate the aggregate national valuation. Conversely, single mean imputation would systematically underestimate variance and distort the relationships between spatial variables \parencite{Burren2018}.

To maintain econometric rigor, the missing acreage parameters were estimated using Multiple Imputation by Chained Equations (MICE). Following the Bayesian framework established by \textcite{Rubin1987}, $m = 5$ parallel, complete datasets were generated. The imputation algorithm utilized Predictive Mean Matching (PMM) configured for 5 convergence iterations (\texttt{maxit = 5}). PMM was selected because it ensures that imputed values are restricted to the range of observed data, preventing the generation of mathematically impossible negative acreages. The predictor matrix for the imputation model included \texttt{course\_id}, \texttt{holes}, \texttt{course\_type}, \texttt{longitude}, \texttt{latitude}, \texttt{state}, and \texttt{state\_land\_value}.

\subsection{Parameter Pooling and Rubin's Rules}
To execute the final analysis, all statistical procedures were performed independently across the $m=5$ imputed datasets. For core parametric modeling, logarithmic regressions (e.g., $\ln(estimated\_value + 1) = \beta_0 + \beta_1 holes + \beta_2 course\_type + \beta_3 county\_type + \epsilon$) were fitted to each dataset. 

The resulting coefficients and standard errors were then statistically merged using Rubin's Rules via the native \texttt{pool()} variance mechanism. This bipartite variance calculation explicitly captures both the within-imputation sampling error and the between-imputation statistical noise injected by the nonresponse \parencite{Burren2018}. Similarly, the aggregate national land value (the scalar estimand) was calculated by summing the total value within each of the five independent datasets and calculating the arithmetic mean of those five sums, ensuring the final aggregate figure accurately reflects the underlying imputation uncertainty.

\newpage

\section{Empirical Results}

This section presents the quantitative findings derived from the spatial valuation model and the multiply imputed datasets. The results quantify the aggregate opportunity cost of golf course land and decompose this value across spatial and operational dimensions.

\subsection{Aggregate Opportunity Cost and Imputation Stability}
The primary objective of this analysis is to quantify the aggregate opportunity cost ($V_{OC}$) of the 16,297 golf courses identified within the United States. Utilizing the hybrid valuation algorithm and pooling the results across the $m=5$ imputed datasets via Rubin's Rules, the aggregate baseline land value is estimated at \$244,922,026,597. 

The standard error of the pooled estimate is \$1.56 billion, yielding a 95\% confidence interval of \$241.86 billion to \$247.98 billion. The narrow confidence interval relative to the aggregate sum indicates high statistical stability within the Predictive Mean Matching (PMM) imputation model, confirming that the between-imputation variance introduced by the 38\% missing acreage data does not compromise the macroeconomic significance of the finding. 

It is critical to note that this \$244.9 billion figure represents a rigorously constrained estimate. Preliminary unbifurcated models, which applied uniform residential values nationally, yielded estimates exceeding \$318 billion. The implementation of the USDA Rural-Urban Continuum Code (RUCC) bifurcation successfully corrected for this upward bias by appropriately assigning lower agricultural counterfactuals ($V_{Ag}$) to inelastic rural markets.

\subsection{Spatial Heterogeneity and Value Concentration}
The aggregate opportunity cost is not distributed uniformly across the spatial landscape. Table \ref{tab:bifurcation} details the bifurcation of the dataset into Urban and Rural classifications.

\begin{table}[h!]
\centering
\caption{Pooled Spatial Bifurcation of Golf Course Opportunity Costs}
\label{tab:bifurcation}
\renewcommand{\arraystretch}{1.2}
\begin{tabular}{l r r r r}
\hline
\textbf{County Type} & \textbf{Observations ($N$)} & \textbf{Total Acreage} & \textbf{Total Estimated Value} & \textbf{Avg. Price Per Acre} \\
\hline
Urban & 11,210 & 1,627,878 & \$241,219,478,404 & \$148,181.30 \\
Rural & 4,877 & 607,060 & \$3,349,241,518 & \$5,517.24 \\
Unknown & 209 & 26,928 & \$353,306,675 & \$13,120.21 \\
\hline
\textbf{Total} & \textbf{16,297} & \textbf{2,261,866} & \textbf{\$244,922,026,597} & \textbf{--} \\
\hline
\end{tabular}
\end{table}

The data reveal an extreme spatial concentration of value. While Urban courses account for 68.8\% of the total facilities and 71.9\% of the total acreage, they represent 98.5\% of the total aggregate opportunity cost. The average price per acre for an Urban course (\$148,181) is nearly 27 times higher than that of a Rural course (\$5,517). 

This finding provides robust empirical confirmation of the theoretical frameworks established by \textcite{Davis2021} and \textcite{Glaeser2005}. The massive opportunity cost of golf course land is a distinctly urban and suburban phenomenon, driven by the ``zoning tax'' and artificial land scarcity in highly regulated metropolitan statistical areas. In rural markets, where land supply is highly elastic, the opportunity cost of maintaining recreational land is macroeconomically negligible.

\subsection{Econometric Determinants of Land Value}
To further analyze the drivers of this opportunity cost, a pooled logarithmic regression was estimated across the $m=5$ datasets. The model regresses the natural log of the estimated land value against physical and operational covariates. The pooled results are presented in Table \ref{tab:regression}.

\begin{table}[h!]
\centering
\caption{Pooled Regression Output: $\ln(Estimated\_Land\_Value + 1)$}
\label{tab:regression}
\renewcommand{\arraystretch}{1.2}
\begin{tabular}{l r r r r}
\hline
\textbf{Term} & \textbf{Estimate ($\beta$)} & \textbf{Standard Error} & \textbf{Statistic ($t$-value)} & \textbf{$p$-value} \\
\hline
(Intercept) & 12.56291 & 0.03719 & 337.79 & $< 0.0001$ \\
holes & 0.04100 & 0.00125 & 32.91 & $< 0.0001$ \\
course\_type: Private & 0.14279 & 0.03979 & 3.59 & $0.00046$ \\
course\_type: Public & -0.23426 & 0.03169 & -7.39 & $< 0.0001$ \\
course\_type: Semi Private & 0.25222 & 0.04465 & 5.65 & $< 0.0001$ \\
county\_type: Unknown & 0.87414 & 0.09702 & 9.01 & $< 0.0001$ \\
county\_type: Urban & 2.69955 & 0.02565 & 105.25 & $< 0.0001$ \\
\hline
\end{tabular}
\end{table}

All covariates in the model are statistically significant at the $p < 0.001$ level. The coefficient for \texttt{holes} ($\beta = 0.041$) indicates that each additional hole increases the underlying land value by approximately 4.1\%, holding other factors constant, which logically scales with the physical footprint required.

Operationally, the model reveals significant heterogeneity. Relative to the baseline, \texttt{Private} courses are associated with a 15.3\% higher underlying land value ($e^{0.14279} - 1 \approx 0.153$), whereas \texttt{Public} courses are associated with a 20.8\% lower land value ($e^{-0.23426} - 1 \approx -0.208$). This aligns precisely with the joint production theory articulated by \textcite{Limehouse2012}; private courses are disproportionately embedded within high-end, master-planned residential communities where the extensive margin value of land is highest.

Finally, the spatial coefficient \texttt{county\_type: Urban} ($\beta = 2.69955$) is the most powerful predictor in the model. Controlling for course size and type, locating a golf course in an urban county increases its opportunity cost by a factor of 14.8 ($e^{2.69955} \approx 14.87$). This mathematically isolates the magnitude of the regulatory scarcity premium ($\tau_z$) that prevents this land from transitioning to its Highest and Best Use.

\newpage

\section{Discussion and Policy Implications}

The empirical results of this study demonstrate that the aggregate opportunity cost of maintaining golf courses in the United States is approximately \$244.9 billion. Crucially, 98.5\% of this foregone value is concentrated in urban and suburban counties. This section contextualizes these findings within broader macroeconomic policy, examines the institutional frictions preventing market optimization, and explores alternative land-use counterfactuals.

\subsection{The Magnitude of Spatial Misallocation}
The \$241.2 billion opportunity cost isolated within urban counties represents a massive spatial misallocation of capital. In highly regulated, inelastic housing markets, the restriction of land supply artificially inflates residential land values ($V_{Res}$) \parencite{Glaeser2005, Davis2021}. By dedicating over 1.6 million acres of prime urban real estate to a single-use recreational amenity, municipalities are implicitly accepting a severe reduction in potential housing density. 

As \textcite{Glaeser2017} note, the inability to deliver housing in high-productivity urban centers restricts labor mobility and depresses aggregate national economic output. Therefore, the \$241.2 billion figure is not merely a localized appraisal metric; it is a quantification of the deadweight loss imposed by maintaining low-density recreational land in areas experiencing acute housing shortages.

\subsection{Institutional Frictions and the Failure of Coasian Bargaining}
If the residential value of this land vastly exceeds its recreational value, standard economic theory suggests that developers should acquire these properties and transition them to their Highest and Best Use (HBU). The failure of the market to execute this transition highlights severe institutional frictions that render Coasian bargaining impossible \parencite{Coase1960}.

The primary friction is the capitalization of the golf course amenity into the surrounding real estate. As \textcite{Cheshire2002} demonstrate, the preservation of open space generates significant gross benefits for adjacent homeowners. Redeveloping the golf course inflicts a direct capital loss on these incumbents, incentivizing fierce political resistance (NIMBYism) to any zoning reform.

Furthermore, even if municipal zoning is amended to permit high-density residential development, private contractual law frequently supersedes public policy. \textcite{Ellickson2022} documents that the vast majority of these developments are governed by restrictive covenants that mandate the preservation of the golf course. Because these covenants require supermajority approval to amend, they suffer from insurmountable collective action problems and status quo bias, effectively freezing the land use in perpetuity. 

To resolve this paralysis, \textcite{Korngold2024} proposes that state legislatures or judiciaries must void these ``stale'' single-family covenants under the doctrine of violating public policy (e.g., the urgent need for housing supply). However, this approach carries significant constitutional risk; voiding a private covenant may be classified as a regulatory or physical taking under the Fifth Amendment, requiring the government to pay ``just compensation'' to the surrounding homeowners \parencite{Korngold2024}. Consequently, the \$244.9 billion opportunity cost remains locked not by a lack of market demand, but by an intractable web of property law and transaction costs.

\subsection{Alternative Counterfactuals: The Renewable Energy Imperative}
Given the severe legal and political frictions associated with residential redevelopment, policymakers must consider alternative HBU counterfactuals. \textcite{Weinand2025} identify utility-scale renewable energy infrastructure as a highly synergistic alternative use for underutilized golf courses. 

The spatial footprint of U.S. golf courses is approximately four times larger than the total land area currently dedicated to utility-scale photovoltaic (PV) solar installations \parencite{Weinand2025}. Repurposing defunct golf courses for solar or wind generation addresses the global ``land squeeze'' without necessarily triggering the same density and traffic concerns that mobilize NIMBY opposition against housing developments. While this study quantified $V_{OC}$ strictly against residential and agricultural baselines, future policy frameworks should integrate $V_{Renewable}$ as a primary counterfactual, particularly in arid urban environments where the ecological cost of maintaining golf course turf exacerbates the economic opportunity cost.

\subsection{Methodological Limitations}
While the hybrid valuation algorithm and MICE protocol ensure a robust baseline estimate, this study acknowledges specific limitations inherent in the data. The FHFA residential land values utilized to calculate $V_{Res}$ are derived from the appraised values of existing, subdivided single-family lots \parencite{Davis2021}. 

In professional appraisal practice, the value of a single 150-acre parcel is not strictly equivalent to the per-acre price of a subdivided lot multiplied by 150. A developer acquiring a bulk parcel requires a discount to account for subdivision costs, infrastructure development (roads, utilities), and the absorption rate of selling the finished lots over time \parencite{AppraisalInstitute2020}. Because the current model applies the marginal per-acre value to the entire bulk acreage, the \$244.9 billion figure should be interpreted as the \textit{gross} theoretical opportunity cost of the land, rather than the immediate liquidation value of the parcels.

\newpage

\section{Conclusion}

This thesis investigated the macroeconomic inefficiency of recreational land allocation in the United States, specifically quantifying the aggregate opportunity cost of maintaining 16,297 golf courses in supply-constrained markets. By integrating geospatial boundary data with granular residential and agricultural land valuations, and by rigorously resolving missing spatial covariates through Multiple Imputation by Chained Equations (MICE), this research established a robust empirical estimate of the foregone economic value of this specific land-use typology.

The empirical results reveal that the aggregate baseline opportunity cost of U.S. golf courses is approximately \$244.9 billion. Crucially, this spatial misallocation is not uniformly distributed. The data demonstrate an extreme concentration of locked land value within metropolitan areas; while urban courses account for roughly 72\% of the total recreational acreage, they represent 98.5\% of the total aggregate opportunity cost. This confirms that the massive financial premium associated with golf course land is fundamentally driven by the regulatory ``zoning tax'' and artificial land scarcity endemic to highly regulated urban and suburban housing markets \parencite{Glaeser2005, Davis2021}. In contrast, the opportunity cost of golf courses in highly elastic, rural markets is macroeconomically negligible.

Despite the residential value of urban golf courses vastly exceeding their capitalized recreational value, the market fails to transition these parcels to their Highest and Best Use (HBU). This thesis concludes that this persistence is not a function of organic market demand, but rather a symptom of severe institutional frictions. A dual layer of public exclusionary zoning and inflexible private restrictive covenants effectively paralyzes land-use transitions, rendering Coasian bargaining impossible and legally barring redevelopment \parencite{Coase1960, Ellickson2022, Korngold2024}. 

Consequently, unlocking this \$244.9 billion in latent land value requires more than market forces; it necessitates aggressive legal and regulatory reform. Future economic research must expand upon this baseline valuation by modeling the specific transaction costs of voiding private covenants and by applying bulk-sale discount (plattage) adjustments to large-acreage parcels. Furthermore, as municipalities grapple with both housing shortages and the global energy transition, future spatial models should integrate utility-scale renewable energy infrastructure as a primary, non-residential HBU counterfactual \parencite{Weinand2025}. Ultimately, the \$244.9 billion opportunity cost quantified in this study serves as a stark financial metric of the deadweight loss imposed by inflexible land-use policies in the modern American economy.

\newpage

\printbibliography[title={References}]

\newpage
\section*{Notes:}

\end{document}